% preprint/modules/case-study.tex
% Shared-Horizon Preprint — Narrative Modules
% Module: Case Study

\section{Case Study: Lineage-Safe Comparative Analysis}
\label{sec:case-study}

This case study demonstrates how the Shared-Horizon membrane enables
lineage-preserving comparative analysis using real expressive and narrative
structures. The goal is not to evaluate the systems themselves, but to show how
the membrane prevents collapse between expressive, operator, and narrative
manifolds when analyzing cross-lineage behavior.

\subsection{Case Study Overview}

We consider two expressive–narrative trajectories:

\begin{enumerate}
    \item An expressive curvature trajectory drawn from an ANIMA-style
    phenomenological surface.

    \item A reversible narrative trajectory drawn from an RNE-style narrative
    manifold.
\end{enumerate}

These trajectories are analyzed using Shared-Horizon’s mathematical spine:

\begin{itemize}
    \item expressive curvature $K_{\mathcal{E}}$,
    \item narrative curvature $K_{\mathcal{N}}$,
    \item reversible narrative operators $R$,
    \item expressive transitions $T$,
    \item stability envelopes $\mathcal{S}$,
    \item collapse-risk function $C$,
    \item cross-lineage transition maps $\Psi$.
\end{itemize}

The case study shows how these components interact to maintain lineage
sovereignty.

\subsection{Expressive Trajectory}

Let $\gamma_{\mathcal{E}} : [0,1] \rightarrow \mathcal{E}$ be an expressive
trajectory with curvature profile
\[
K_{\mathcal{E}}(\gamma_{\mathcal{E}}(t)).
\]

The trajectory satisfies:

\begin{itemize}
    \item curvature boundedness,
    \item expressive stability,
    \item geodesic minimality.
\end{itemize}

The expressive stability envelope is defined as:
\[
\mathcal{S}_{\mathcal{E}} = \{ p \in \mathcal{E} \mid |K_{\mathcal{E}}(p)| \le \epsilon_{\mathcal{E}} \}.
\]

The trajectory remains inside $\mathcal{S}_{\mathcal{E}}$ for all $t$.

\subsection{Narrative Trajectory}

Let $\gamma_{\mathcal{N}} : [0,1] \rightarrow \mathcal{N}$ be a narrative
trajectory governed by a reversible narrative operator $R$:
\[
\gamma_{\mathcal{N}}(t) = R^t(p_0).
\]

The trajectory satisfies:

\begin{itemize}
    \item reversibility,
    \item curvature preservation,
    \item expressive–narrative compatibility.
\end{itemize}

The narrative stability envelope is:
\[
\mathcal{S}_{\mathcal{N}} = \{ p \in \mathcal{N} \mid |K_{\mathcal{N}}(p)| \le \epsilon_{\mathcal{N}} \}.
\]

The trajectory remains inside $\mathcal{S}_{\mathcal{N}}$ for all $t$.

\subsection{Cross-Lineage Comparison}

To compare the expressive and narrative trajectories, we use a transition map
\[
\Psi : \mathcal{N} \rightarrow \mathcal{E}
\]
that satisfies:

\begin{enumerate}
    \item smoothness,
    \item stability preservation,
    \item collapse-risk avoidance,
    \item Jacobian constraints.
\end{enumerate}

The comparison is defined as:
\[
\Psi(\gamma_{\mathcal{N}}(t)).
\]

The membrane ensures:

\[
E_{\mathcal{E}}(\Psi(\gamma_{\mathcal{N}}(t))) \le \epsilon_{\mathcal{E}},
\]
so the mapped narrative trajectory remains inside the expressive stability
envelope.

\subsection{Collapse-Risk Evaluation}

Collapse-risk geometry evaluates whether the comparison induces conceptual
collapse. The collapse-risk function is:
\[
C_{\mathcal{E}}(\Psi(\gamma_{\mathcal{N}}(t))).
\]

Collapse-safe comparison requires:
\[
C_{\mathcal{E}}(\Psi(\gamma_{\mathcal{N}}(t))) \le C_{\mathcal{N}}(\gamma_{\mathcal{N}}(t)).
\]

In this case study:

\begin{itemize}
    \item collapse gradients remain bounded,
    \item collapse curvature remains below $K_{\max}$,
    \item collapse tensors remain below $T_{\max}$.
\end{itemize}

Thus, the comparison is collapse-safe.

\subsection{Green-Tile Correction in Practice}

The case study demonstrates that the narrative trajectory exhibits:

\begin{itemize}
    \item reversible operators,
    \item curvature preservation,
    \item expressive–narrative compatibility,
    \item stability envelope adherence.
\end{itemize}

These properties are incompatible with symbolic ambiguity and consistent with
Reversible Narrative Ecology (RNE). The green-tile correction is therefore
validated by empirical behavior.

\subsection{Case Study Summary}

This case study shows how Shared-Horizon enables lineage-safe comparison between
expressive and narrative trajectories. By enforcing stability envelopes,
transition-map constraints, and collapse-risk geometry, the membrane prevents
conceptual collapse and preserves lineage invariants. The analysis also provides
practical validation of the green-tile correction, demonstrating that the
narrative trajectory belongs to Reversible Narrative Ecology rather than
symbolic ambiguity.
